
\section{Sharpness and fluctuation lower bounds}

We continue with $\rho(x)=|x|$ and put $u_x=x/|x|$ for
$x\ne0$, $u_0=0$.
Write $r=r_n$, $L=\log n$, $R_-=\Rin(n)$, and $R_+=\Rout(n)$.
All estimates below allow arbitrary first-cookie covariances.
We use the already proved quantitative shape estimates on an event
$\Omega_n$, for a fixed constant $K_0<\infty$, with $\mathbb P(\Omega_n^c)\leq Cn^{-10}$:
\begin{align}
 R_+&\leq K_0r,&
 \sup_x|\ell_n(x)-2d\kappa(r-|x|)_+|
      &\leq Cr^{1/2}L^{3/2},\label{eq:lower-input-profile}\\
 |R_--r|+r^{1-d}|D_n\mathbin\triangle B_r|
      &\leq C h_d,&
 h_d&=\begin{cases}\sqrt{rL},&d=2,\\L,&d\geq3.
                  \end{cases}\label{eq:lower-input-volume}
\end{align}
The discrete-to-continuum potential estimate and Dynkin's identity
then give, uniformly for $|y|\leq r/2$,
\begin{equation}\label{eq:lower-bulk-representation}
 \ell_n(y)-2d\kappa(r-|y|)=-M_n^y+O(h_d+L)
 \quad\hbox{on }\Omega_n.
\end{equation}
Indeed, $E=D_n\setminus B_{R_-}$ lies a distance of order $r$
from such $y$, so its drift potential has absolute value at most
$Cr^{1-d}|E|\leq Ch_d$; the potential of $B_{R_-}$ is the cone
of height $2d\kappa R_-$, and the remaining error is $O(L)$.
\subsection{One exponential-martingale lemma}

\begin{lemma}\label{lem:lower-exp}
Fix $0<c_0\leq C_0$.  Let $S^1,\ldots,S^m$ be real
martingales starting at zero with increments bounded by $b$.
Suppose an event $\Omega$, of probability at least $1-\alpha$,
satisfies $c_0\leq\langle S^i\rangle_n\leq C_0$ for every $i$.
There are $c,C>0$, depending only on $c_0,C_0$, such that if
$\beta\geq C$ and $\beta b\leq c$, then
\[
 \mathbb P(\pm S_n^i\geq c\beta)
       \geq\tfrac14e^{-C\beta^2}-\alpha.
\]
If additionally
$|\langle S^i,S^j\rangle_n|\leq\delta$ on $\Omega$ for $i\ne j$
and $\beta^2\delta+\beta^3b\leq1$, then
\[
 \mathbb P\{\max_i S_n^i<c\beta\}
 \leq C\left(m^{-1}e^{C\beta^2}+\beta^2\delta+
                    \beta^3b+e^{C\beta^2}\sqrt\alpha\right).
\]
\end{lemma}
\begin{proof}
Stop all martingales just before any bracket would exceed $2C_0$;
this changes none of them on $\Omega$.  For a martingale $S$ whose
increments are at most $b_1$ and whose bracket is bounded by $8C_0$,
put
\[
 K_S(s)=\sum_{j<n}\log\mathbb E(e^{s\Delta S_j}\mid\mathcal F_j),
 \qquad Z_S(s)=e^{sS_n-K_S(s)}.
\]
Conditional Taylor expansion, for $|s|b_1\leq1/2$, gives
\[
 K_S(s)=\tfrac12s^2\langle S\rangle_n+O(|s|^3b_1),
 \quad \mathbb E Z_S(s)=1,
 \quad \mathbb E Z_S(s)^2\leq e^{Cs^2}.
\]
For the last bound use
$Z_S(s)^2=Z_S(2s)e^{K_S(2s)-2K_S(s)}$.
Apply these formulas to $S^i$ and $S^i+S^j$, with $s=\beta$,
and abbreviate $Z_i=Z_{S^i}(\beta)$.
On $\Omega$, $Z_i\geq1/2$ implies $S_n^i\geq c\beta$,
since $K_{S^i}(\beta)\geq c_0\beta^2/4$.
Paley--Zygmund therefore gives the first assertion; use $-S^i$
for the opposite sign.
For $i\ne j$,
\[
 Z_iZ_j=Z_{S^i+S^j}(\beta)
       \exp\{\beta^2\langle S^i,S^j\rangle_n+O(\beta^3b)\}.
\]
On $\Omega$ the exponential is at most
$1+C(\beta^2\delta+\beta^3b)$; on its complement apply
Cauchy--Schwarz and the second-moment bound.  Consequently
\[
 \mathbb E Z_iZ_j\leq1+C(\beta^2\delta+\beta^3b)
                              +e^{C\beta^2}\sqrt\alpha.
\]
The average $\overline Z=m^{-1}\sum_i Z_i$ has mean one and
variance bounded by the right side of the asserted probability
estimate.  If $\overline Z\geq1/2$ on $\Omega$, some
$S_n^i\geq c\beta$.  Chebyshev, with the additional error $\alpha$,
finishes the proof.
\end{proof}

\subsection{Planar radii and width}

Let $d=2$.  Define the weighted moment and vector martingales
\[
 S_n=\sum_{x\in A_n}\lambda_x|x|,\qquad
 Q_n=|X_n|^2-n+2\kappa S_n,\qquad
 Z_n=X_n+\kappa\sum_{x\in A_n}\lambda_xu_x.
\]
On $\{R_+\leq K_0r\}$ the defect assumption gives
\begin{equation}\label{eq:lower-defect-scores}
 \left|S_n-\sum_{x\in A_n}|x|\right|\leq Cr^2,
 \qquad
 \left|\sum_{x\in A_n}(\lambda_x-1)u_x\right|\leq Cr.
\end{equation}
The first-cookie covariance corrections contribute $O(r^4)$ to
$\langle Q\rangle_n$ and $O(r^2)$ to each bracket of $Z$.
Thus \eqref{eq:lower-input-profile} gives, on $\Omega_n$,
\begin{equation}\label{eq:lower-QZ-brackets}
 \langle Q\rangle_n=(4\pi\kappa/5+o(1))r^5,
 \qquad
 \langle v\cdot Z\rangle_n=|v|^2n/2+O(|v|^2r^2).
\end{equation}
These formulas also hold almost surely with relative error $o(1)$
by the qualitative shape theorem.

The unweighted moment radius
\[
 T_n=\left(\frac3{2\pi}\sum_{x\in A_n}|x|\right)^{1/3}
\]
satisfies $R_-\leq T_n+C$ and $R_+\geq T_n-C$ by comparison
with unit cells.  The weighted moment radius differs from $T_n$
by $O(1)$ on $\Omega_n$ and eventually almost surely, by
\eqref{eq:lower-defect-scores} and the shape theorem.
The general moment LIL therefore proves
\begin{equation}\label{eq:lower-radii-lil}
 \limsup_n\frac{r-R_-}{\sqrt{r\log\log n}},\quad
 \limsup_n\frac{R_+-r}{\sqrt{r\log\log n}}
 \ \geq\frac1{\sqrt{10\pi\kappa}}\quad\hbox{almost surely}.
\end{equation}
To obtain polynomial lower tails, stop $Q$ on exiting $B_{K_0r}$
and apply Lemma~\ref{lem:lower-exp} to $Q/r^{5/2}$.
Its increments are bounded by $4K_0r^{-3/2}$, and its bracket lies
between positive constants on $\Omega_n$ by
\eqref{eq:lower-QZ-brackets}.
Choose $\beta=a_p\sqrt{\log n}$, with $a_p>0$ small enough that
$Ca_p^2<\min(p,10)/2$.  Each sign then has probability exceeding
$n^{-p}$ after intersecting with $\Omega_n$ and subtracting its
failure probability once more, for large $n$.
Since $|X_n|^2$ and the defect in \eqref{eq:lower-defect-scores}
are $O(r^2)=o(r^{5/2}\sqrt{\log n})$, the identity for $Q$ and
factoring $T_n^3-r^3$ yield, with a suitable $c_p>0$,
\begin{equation}\label{eq:lower-radii-tails}
 \mathbb P\{r-R_-\geq c_p\sqrt{r\log n}\}\geq n^{-p},
 \qquad
 \mathbb P\{R_+-r\geq c_p\sqrt{r\log n}\}\geq n^{-p}.
\end{equation}

For the width set $V_n=\sum_{x\in A_n}u_x$.
The sum over the lattice disk $\{|x|<R_-\}$ vanishes by symmetry.
When $R_-,R_+\asymp r$, integration of the positive part of any
directional component over the remaining annulus gives
\begin{equation}\label{eq:lower-width-geometry}
 |V_n|\leq R_+^2-R_-^2+Cr,
 \qquad |Z_n|\leq\kappa(R_+^2-R_-^2)+Cr.
\end{equation}
Indeed, $\int_0^{2\pi}(\cos\theta)_+\,d\theta=2$; replacing
each unit vector by its values in the corresponding cell costs
$O(r)$, as does the weighted correction.
The martingale $e_1\cdot Z$ has increments at most two.
Apply Lemma~\ref{lem:lower-exp} to $(e_1\cdot Z)/\sqrt n$,
using its bracket bounds on $\Omega_n$ from
\eqref{eq:lower-QZ-brackets}.  Equation
\eqref{eq:lower-width-geometry} and $n=(4\pi\kappa/3)r^3$ give
\begin{equation}\label{eq:lower-width-tail}
 \mathbb P\{R_+-R_-\geq c_p\sqrt{r\log n}\}\geq n^{-p}.
\end{equation}
The same bracket limit, the bounded-increment martingale LIL,
and $R_-+R_+\sim2r$ give
\begin{equation}\label{eq:lower-width-lil}
 \limsup_n\frac{R_+-R_-}{\sqrt{r\log\log n}}
       \geq\sqrt{\frac\pi{3\kappa}}\quad\hbox{almost surely}.
\end{equation}
Finally the vector martingale CLT yields
$Z_n/\sqrt n\Longrightarrow N(0,I_2/2)$.
On $\{R_+-R_-\leq a\sqrt r\}$,
\eqref{eq:lower-width-geometry} gives
$|Z_n|/\sqrt n\leq\sqrt{3\kappa/\pi}\,a+o(1)$ almost surely.
Since $\mathbb P\{|N(0,I_2/2)|\leq t\}=1-e^{-t^2}$,
\begin{equation}\label{eq:lower-width-small}
 \limsup_n\mathbb P\{R_+-R_-\leq a\sqrt r\}
       \leq1-e^{-3\kappa a^2/\pi}\qquad(a\geq0).
\end{equation}

\subsection{The maximum of the bulk local times}

Let $m=\lfloor r^{1/8}\rfloor$,
$y_i=i\lceil r^{1/4}\rceil e_1$, and put
\[
 f_i=\begin{cases}g_{y_i+ke_1}-g_{y_i},\quad k=\lfloor r^{1/8}\rfloor,
                         &d=2,\\g_{y_i},&d\geq3,
       \end{cases}
 \qquad
 s=\begin{cases}\sqrt{r\log r},&d=2,\\\sqrt r,&d\geq3,
      \end{cases}
 \qquad S^i=-M^{f_i}/s.
\]
All the poles lie in $B_{\sqrt r}$ for large $n$.
We verify, uniformly on $\Omega_n$,
\begin{equation}\label{eq:lower-separated}
 c_0\leq\langle S^i\rangle_n\leq C_0,
 \qquad |\langle S^i,S^j\rangle_n|\leq Cr^{-1/4}\quad(i\ne j).
\end{equation}
This calculation uses only the simple-random-walk covariance kernel
$\Gamma$ and the cone estimate; no independence is claimed.
Write $W_{ij}=\sum_x|\Gamma(f_i,f_j)(x)|$ and
$T_{ij}=\sum_x\min(|x|,r)|\Gamma(f_i,f_j)(x)|$.
The Green-kernel gradient bound and $2ab\leq a^2+b^2$ give
\[
 (W_{ij},T_{ij})\leq
 \begin{cases}(C\log r,C\sqrt r\log r),&d=2,\\
                  (C,C(\sqrt r+\log r)),&d\geq3.
       \end{cases}
\]
For completeness, in $d\geq3$ the squared increments are bounded
by $C(1+|x-y_i|)^{2-2d}$; its sum is finite and its first truncated
moment is $O(\log r)$ in $d=3$ and $O(1)$ in larger dimensions.
In $d=2$, an increment of $f_i$ is bounded by the sum of reciprocal
distances to its two poles within distance $3k$, and by
$Ck(1+|x-y_i|)^{-2}$ outside.  Squaring and summing gives
$O(\log r)$ and a first moment $O(k\log r)$, respectively.
Translation of the poles costs at most $C\sqrt r$ times the
unweighted sums.  The finitely supported Laplacian terms in
$\Gamma$ satisfy the same bounds.

Replacing first-cookie covariances by $\Gamma$ costs $O(1)$ in
$d\geq3$ and $O(\log r)$ in $d=2$.
Replacing $\ell_n$ by its cone and then by the constant $2d\kappa r$
costs at most $Cr^{1/2}L^{3/2}W_{ij}+2d\kappa T_{ij}$.
After division by $s^2$ all these errors are
$O(r^{-1/2}L^{3/2})=o(r^{-1/4})$.
Thus the bracket differs by this error from
\[
 \Sigma_{ij}=\frac{2d\kappa r}{s^2}
                         \sum_x\Gamma(f_i,f_j)(x).
\]
For $d\geq3$, summation by parts gives
$\Sigma_{ij}=2d\kappa(2G(y_i-y_j)-\mathbf1_{\{i=j\}})$.
Its diagonal is positive, and its off-diagonal entries are
$O(r^{(2-d)/4})$.
For $d=2$, writing $z=y_i-y_j$, the same identity gives
\[
 \sum_x\Gamma(f_i,f_j)(x)
   =2\{g(z+ke_1)+g(z-ke_1)-2g(z)\}-2\mathbf1_{\{i=j\}}.
\]
The diagonal of $\Sigma$ tends to $4\kappa/\pi$, because
$g(ke_1)=(2/\pi)\log k+O(1)$.
For $i\ne j$, the displayed second difference is
$O(k^2/|z|^2)$, giving an off-diagonal bound $O(r^{-1/4})$.
This proves \eqref{eq:lower-separated}.

The normalized martingale increments are bounded by $Cr^{-1/2}$.
Apply Lemma~\ref{lem:lower-exp} with
$\beta=a\sqrt{\log r}$, where $a>0$ is so small that
$Ca^2<1/16$.  Since $m\asymp r^{1/8}$,
$\delta=Cr^{-1/4}$, and $\alpha=O(n^{-10})$, the probability
that all $S_n^i<c\sqrt{\log r}$ is at most $Cr^{-1/16}$.
The error in \eqref{eq:lower-bulk-representation} is
$o(s\sqrt{\log r})$.  In $d\geq3$ a large $S_n^i$ gives a large
positive local-time deviation at $y_i$; in $d=2$ it gives a large
difference of the deviations at its two poles, so one has large
absolute value.  Consequently, for some $c>0$,
\begin{equation}\label{eq:lower-bulk-max}
 \mathbb P\left\{\max_{|y|\leq\sqrt r}
 |\ell_n(y)-2d\kappa(r-|y|)|\geq
 c\begin{cases}\sqrt r\log n,&d=2,\\\sqrt{r\log n},&d\geq3
   \end{cases}\right\}\geq1-n^{-c}
\end{equation}
for all sufficiently large $n$.

\subsection{Logarithmic radii under ellipticity}

For this paragraph assume additionally $q_x(e)\geq p_0>0$
for every $x,e$; later departures already have probability $1/(2d)$.
There are $c,C,h_0>0$ such that, for $h_0\leq h\leq r/8$,
\begin{align}
 \mathbb P\{R_-\geq r-h,\ R_+\leq r+h\}
   &\leq e^{-cr^{d-1}e^{-Ch}},\label{eq:lower-log-joint}\\
 \mathbb P\{R_+\leq r+h\}
   &\leq e^{-cr^{d-1}e^{-Ch}}+e^{-cr^{d-3}h^2}.
       \label{eq:lower-log-outer}
\end{align}
To prove these estimates, from any $x$ with $|x|>r-2h$ take
$T=\lceil4h\sqrt d\rceil$ steps in the signed coordinate of
largest $|x_i|$.  The Euclidean distance increases at speed at least
$1/\sqrt d$, so this path exits $B_{r+h}$; its conditional
probability is at least $p_0^T$ after decreasing $p_0$ if necessary.
Among $k$ distinct arrivals at sites of any fixed
$\Lambda\subset\{|x|>r-2h\}$ by time $n$, at least
$\lfloor(k-T)/(T+1)\rfloor$ disjoint such trials can be started
and completed by time $n$.  Successive conditioning bounds their
joint failure probability by
$(1-p_0^T)^{\lfloor(k-T)/(T+1)\rfloor}$.
If $R_-\geq r-h$, the annulus
$r-2h<|x|<r-h-1$ supplies $k\geq cr^{d-1}h$ sites, proving
\eqref{eq:lower-log-joint}.

For \eqref{eq:lower-log-outer}, let
$\Lambda=\{r-2h<|x|\leq r+h\}$.
The same trials apply if $|A_n\cap\Lambda|\geq c_1r^{d-1}h$.
Otherwise, on $R_+\leq r+h$, cell comparison and the defect bound
give
\[
 \sum_{x\in A_n}\lambda_x|x|
 \leq\frac{n}{2\kappa}-c_2r^dh
                     +2r|A_n\cap\Lambda|+Cr^d.
\]
Choose $c_1$ small and then $h_0$ large.  The identity
$Q_n=|X_n|^2-n+2\kappa\sum\lambda_x|x|$
forces $Q_n\leq-cr^dh$; the endpoint $|X_n|^2$ is absorbed
by $h_0$ in $d=2$ and by large $r$ in higher dimensions.
Stop $Q$ on leaving $B_{r+h}$.  Its increments are bounded by
$4(r+h)$, so Azuma's inequality gives
$\mathbb P\{Q_n\leq-cr^dh,\ R_+\leq r+h\}
\leq e^{-cr^{d-3}h^2}$, proving the second estimate.
Taking $h=c\log n$ with $c>0$ sufficiently small and applying
Borel--Cantelli gives, almost surely for all sufficiently large $n$,
\[
 \max(r-R_-,R_+-r)>c\log n;
 \qquad R_+-r>c\log n\quad\hbox{if }d\geq3.
\]
The ellipticity assumption is used only in the prescribed-path
trials and is not needed for any preceding lower bound.
